% TOP_PROBLEM: {"id":30007542,"problem_number":"LOCAL-30007542","title":"Research productivity versus commercialization","category_id":21,"category":{"id":21,"name":"economics","display_name":"Economics","description":"Open economic theory statements, published research agendas, and proposed scoped specifications; question provenance and historical priority are qualified per record.","slug":"economics","order_index":21,"created_at":"2026-10-08T00:00:00Z"},"status":"research_question","external_url":null,"legacy_ids":[],"published":false,"statement":"Identify declining research productivity separately from patent measurement and implementation barriers in a latent-idea growth model.","economics":{"original_id":31,"field":"Growth and productivity","kind":"identification","formalization_basis":"adapted_research_question","source_status":"research_agenda","priority_status":"earliest_found","priority_note":"This is the earliest explicit relevant statement verified in this audit, not a claim of original priority; older literature and earlier versions may contain antecedents.","earliest_verified_reference":{"authors":["Teresa C. Fort","Nathan Goldschlag","Jack Liang","Peter K. Schott","Nikolas Zolas"],"title":"Growth is Getting Harder to Find, Not Ideas","year":2025,"url":"https://www.census.gov/library/working-papers/2025/adrm/CES-WP-25-21.html","doi":null,"locator":"Abstract, concluding paragraph, Census CES-25-21, April 2025","evidence_summary":"The paper challenges a universal decline in researcher productivity using patent evidence and explicitly invites renewed empirical and theoretical attention to how ideas affect economic growth. This is an agenda, not a canonical mathematical conjecture."},"current_reference":{"authors":["Teresa C. Fort","Nathan Goldschlag","Jack Liang","Peter K. Schott","Nikolas Zolas"],"title":"Growth is Getting Harder to Find, Not Ideas","year":2025,"url":"https://www.census.gov/library/working-papers/2025/adrm/CES-WP-25-21.html","doi":null,"locator":"Abstract, concluding paragraph, Census CES-25-21, April 2025","evidence_summary":"The paper challenges a universal decline in researcher productivity using patent evidence and explicitly invites renewed empirical and theoretical attention to how ideas affect economic growth. This is an agenda, not a canonical mathematical conjecture."},"formalization_note":"The paper explicitly calls for renewed work on ideas’ growth effects and challenges a universal research-productivity decline. This latent-variable decomposition is adapted."},"statusQualification":"Published question or research agenda verified at the cited locator. The mathematical specification is adapted for this catalogue and includes additional assumptions; source publication does not establish first-ever priority or certify this exact model remains unsolved. The paper explicitly calls for renewed work on ideas’ growth effects and challenges a universal research-productivity decline. This latent-variable decomposition is adapted.","statusReviewedAt":"2026-10-08"}
% FIRST_POSTED: 2026-10-08
% ENTRY_KIND: statement_only
% !TeX program = lualatex
\documentclass[11pt]{article}
\usepackage{fontspec}
\usepackage[a4paper,margin=25mm]{geometry}
\usepackage{amsmath,amssymb,mathtools}
\usepackage{xurl}
\usepackage[hidelinks]{hyperref}
\newcommand{\sref}[2]{\hyperlink{source:#1:#2}{[S#2]}}
\setlength{\emergencystretch}{3em}
\begin{document}
% problemId:30007542
\section*{Research productivity versus commercialization}\label{30007542}
\subsection{Definitions and mathematical statement}
Fix industries \(j\) and years \(t\). Let \(R_{jt}\) measure research input, \(I_{jt}\) latent economically useful new ideas, \(P_{jt}\) observed patents adjusted by a predeclared quality measure, and \(A_{jt}\) measured productivity. Specify a research relation \(I_{jt}=F_\theta(R_{jt},A_{jt},v_{jt})\), a patent measurement relation \(P_{jt}=M_\psi(I_{jt},z_{jt})+\epsilon_{jt}\), and an adoption relation \(A_{j,t+h}=D_\gamma(I_{\le t},K_{jt},b_{jt},h)+\nu_{j,t+h}\). Here \(v\) denotes research opportunities, \(z\) patenting incentives, \(K\) complementary capital, and \(b\) commercialization barriers.
Identify changes in research productivity \(\partial F_\theta/\partial R\) separately from changes in patent measurement and adoption. The target is a decomposition of a declared productivity-growth slowdown into research-output and implementation components, with interactions allocated by an explicitly stated counterfactual convention. Use variation affecting research effort and adoption costs differently, explaining exclusion restrictions and spillovers across firms. Test whether competing parameterizations fit both patent-quality distributions and independently measured product or process improvements. Report identified sets when latent idea quality prevents point identification. Patent counts alone cannot establish that ideas become uniformly harder to find, while rising patents alone cannot establish increasing economically relevant research productivity.

\par\textbf{Formulation and scope.} The paper explicitly calls for renewed work on ideas’ growth effects and challenges a universal research-productivity decline. This latent-variable decomposition is adapted.

\par\textbf{Open-question provenance.} An explicit question or research agenda was verified at the source locator. This does not by itself certify that every later version remains unresolved \sref{30007542}{1}.

\par\textbf{Historical priority.} This is the earliest explicit relevant statement verified in this audit, not a claim of original priority; older literature and earlier versions may contain antecedents.

\subsection{Short English statement}
Identify declining research productivity separately from patent measurement and implementation barriers in a latent-idea growth model.

\subsection{Sources}
\begin{itemize}\raggedright
\item[\textnormal{[S1]}]\hypertarget{source:30007542:1}{} Teresa C. Fort, Nathan Goldschlag, Jack Liang, Peter K. Schott, Nikolas Zolas. 2025. Growth is Getting Harder to Find, Not Ideas. Abstract, concluding paragraph, Census CES-25-21, April 2025.
\url{https://www.census.gov/library/working-papers/2025/adrm/CES-WP-25-21.html}.
{\small\textit{Earliest verified question/agenda reference; historical priority qualified below. The paper challenges a universal decline in researcher productivity using patent evidence and explicitly invites renewed empirical and theoretical attention to how ideas affect economic growth. This is an agenda, not a canonical mathematical conjecture.}}
\end{itemize}
\end{document}
